\author{OTS}
\documentclass[14pt]{../../inc/mybib}

\setincpath{../../inc/}

\usepackage{bible_style}
\usepackage{bibeltext}
\usepackage{markierungen}
\graphicspath{{../../assets/images/}}
\usepackage{header}

\newcommand{\Name}{Fredy}

% ensure scrlayer-scrpage has sufficient footheight
\setlength{\footheight}{10.4pt}


\begin{document}
\begin{bibelbox}{SCHL}{Kol}{3:15-17}
    Lasst das Wort des Christus reichlich in euch wohnen in aller Weisheit; lehrt und ermahnt einander und singt mit Psalmen und Lobgesängen und geistlichen Liedern dem Herrn lieblich in eurem Herzen. Und was immer ihr tut in Wort und Werk, das tut alles im Namen des Herrn Jesus und dankt Gott, dem Vater, durch ihn.
\end{bibelbox}
\section{Begrüssung}
Ich möchte auch euch alle recht herzlich zum Gottesdienst begrüssen. Schön, dass ihr gekommen seid, um unserem \herr N zu danken, ihn zu loben und zu preisen.

Auch dich \Name{} will ich herzlich begrüssen. Schön, dass du wieder bei uns bist und mit uns Gottesdienst feierst.
\begin{itemize}
    \item[] \beten[b]{Beten zur Begrüssung}
    \item[] Und anschliessend singen wir zusammen das Lied
    \item[] \lied[b]{Heilig, heilig, heilig, Gott, Du bist heilig}
\end{itemize}

\section{Informationen}
\begin{itemize}
    \item \info[b,p]{Bibel und Gebetsabend:} Do, 23.April.2026 20:00 Uhr Bibel und Gebetsabend mit Markus 6,45-
    \item \info[b,p]{Nächster Gottesdienst:} So, 26.April.2026 14:45 Uhr dieses Mal Qumran und die Jesaja-Rolle

    \item \info[b,p]{Kollekte:} Die Kollekte wird hinten in der grünen Box gesammelt und wird vollumfänglich für den Bau dieser Gemeinde eingesetzt.
\end{itemize}

\section{ Input }
\begin{spacing}{1.5}
    \begin{block}[Gott hat die Kontrolle? Wer aber ist Gott?]
        Ein halbes Jahr lang habe ich euch die Heilsgeschichte Gottes dargelegt. Ich habe oft betont, wie souverän Gott ist, dass er die Kontrolle über uns und über die ganze Welt hat. Ich habe versucht, dies mit der Bibel zu belegen. Aber wer ist eigentlich dieser Gott? Niemand hat ihn je gesehen, trotzdem soll es ihn geben.

        Es gibt viele Vorstellungen, wer Gott ist und wie Gott ist. 
        \begin{itemize}
            \item Für die Feministen ist Gott eine Frau, 
            \item für die Esoteriker ist er eine Energie,
            \item für die Pantheisten ist Gott alles und es gibt keinen persönlichen Gott,
            \item für die Polytheisten gibt es viele Götter,
            \item für die Atheisten ist es ein Urknall.
        \end{itemize}       

        In der Geschichte hat jedes Volk sich einen oder mehrere Götter geschaffen. Oft sind diese Götter mit der Natur verbunden und sollen vor Unheil schützen und Wohlstand bringen.

        Was auffällt ist, dass die meisten Naturvölker eine sehr brutale Moralvorstellung haben. Im Hinduismus war das Töten der Mädchen und das Verbrennen von Witwen normal. Im Islam ist noch heute die Sklavenhaltung und die Erniedrigung der Frau legal. In vielen Naturvölkern waren Genitalverstümmelung, Menschenopfer, Kannibalismus, usw. normal.

        Aber auch unter den Christen gibt es viele verschiedene Vorstellungen, so z.\,B. dass es einen Gott im Alten Testament gibt, der Zornige, der die Massenmorde befohlen hat und dann der Gott aus dem Neuen Testament, der Liebe, der alle Menschen rettet.

        Also was nun? Gibt es \betonung[b,p]{den} einen Gott? Ja! Ich bin überzeugt, dass es \betonung[b,p]{den einen wahren Gott} gibt. Wenn ich alle anderen Götter so ansehe, sind sie alle sehr beschränkt und menschlich. Zum Beispiel Allah, der sich ein Paradies erschaffen hat, mit etlichen Jungfrauen für die guten Muslime. Die Hindus haben 330 Millionen Götter. Der Buddhismus hat gar keinen Gott.

        Keiner dieser Religionen hat einen Gott wie der aus der Bibel. Nirgends in der Bibel findet ihr, woher Gott kommt oder wie er entstanden ist. In anderen Mythen gibt es immer schon etwas und daraus wurde dann ein Gott kreiert. Auch der Atheist sagt, dass am Anfang Energie da war. In unserer Bibel heisst einfach:
        \begin{bibelbox}{SCHL}{IMos}{1:1}
            Im Anfang schuf Gott die Himmel und die Erde.
        \end{bibelbox}
        Das ist alles. Gott ist schon da und er erschuf. Das Gleiche lesen wir auch im Johannes-Evangelium:
        \begin{bibelbox}{SCHL}{Joh}{1:1-3}
            Im Anfang war das Wort, und das Wort war bei Gott, und das Wort war Gott. Dieses war im Anfang bei Gott. Alle Dinge sind durch dasselbe entstanden: und ohne dasselbe ist auch nicht eines entstanden, was entstanden ist.
        \end{bibelbox}
        Ich habe jetzt aus unserer Bibel zitiert. Weil, \betonung[b,p]{ohne} diese Bibel wüssten wir nichts über Gott. Es gibt keine andere Quelle, die uns über Gott informiert. Ich weiss, wenn jemand dieses Buch als Märchenbuch abtut, dann glaubt er auch nicht, was darin steht. Aber wir hier, die wir uns Christen nennen, glauben, dass die Bibel Gottes Wort ist. Und wenn das so ist, dann sollten wir auch glauben, was Gott über sich selbst sagt.

        Was aber sagt Gott über sich selbst?
        \begin{itemize}
            \item \verbN[b,p]{Er ist der eine wahre Gott. Es gibt keinen anderen neben ihm}.
            \begin{bibelbox}{SCHL}{VMos}{6:4}
                Höre, Israel: Der HERR, unser Gott, der HERR allein.
            \end{bibelbox}
            \begin{bibelbox}{SCHL}{Jes}{45:5}
                Ich bin der HERR, und sonst ist keiner; denn außer mir gibt es keinen Gott.
            \end{bibelbox}
            \begin{bibelbox}{SCHL}{IKor}{8:4}
                Was nun das Essen der Götzenopfer betrifft, so wissen wir, dass ein Götze in der Welt nichts ist, und dass es keinen Gott gibt ausser dem Einen.
            \end{bibelbox}
            \item \verbN[b,p]{Er ist ewig und aus sich selbst heraus. Er ist nicht geworden, sondern ist.} 
            \begin{bibelbox}{SCHL}{Ps}{90:2}
                Ehe die Berge wurden und du die Erde und den Erdkreis hervorbrachtest, ja, von Ewigkeit zu Ewigkeit bist du Gott!
            \end{bibelbox}
            \begin{bibelbox}{SCHL}{IIMos}{3:14}
                Gott sprach zu Mose: \flqq Ich bin der ich bin\frqq. Und er sprach: \flqq So sollst du zu den Kindern Israels sagen: \flq Ich bin\frq{} der hat mich zu euch gesandt.\frqq
            \end{bibelbox}
            \begin{bibelbox}{SCHL}{Offb}{1:8}
                Ich bin das A und das O, der Anfang und das Ende, spricht der Herr, der ist und der war und der kommt, der Allmächtige.
            \end{bibelbox}
            \item \verbN[b,p]{Er ist der Schöpfer von Himmel und Erde und der Ursprung allen Lebens.} 1Mo 1,1; Neh 9,6; Apg 17,24-25
            \begin{bibelbox}{SCHL}{IMos}{1:1}
                Im Anfang schuf Gott die Himmel und die Erde.
            \end{bibelbox}
            \begin{bibelbox}{SCHL}{Neh}{9:6}
                Du allein bist der HERR. Du hast Himmel und Erde gemacht, das Meer und alles, was darin ist, und du erhältst alles am Leben.
            \end{bibelbox}
            \begin{bibelbox}{SCHL}{Apg}{17:24-25}
                Der Gott, der die Welt gemacht hat und alles, was darin ist, er, der Herr des Himmels und der Erde, wohnt nicht in Tempeln, die von Menschenhand gemacht sind, und wird auch nicht von Menschenhänden bedient, als ob er etwas nötig hätte, da er selbst allen Leben und Atem und alles gibt.
            \end{bibelbox}
            \item \verbN[b,p]{Er ist heilig, rein und moralisch vollkommen.} Jes 6,3; 3Mo 11,44; 5Mo 32,4
            \begin{bibelbox}{SCHL}{Jes}{6:3}
                Und sie riefen einander zu und sprachen: Heilig, heilig, heilig ist der HERR Zebaoth, alle Lande sind seiner Ehre voll!
            \end{bibelbox}          
            \begin{bibelbox}{SCHL}{Neh}{9:6}
                Du allein bist der HERR. Du hast Himmel und Erde gemacht, das Meer und alles, was darin ist, und du erhältst alles am Leben.
            \end{bibelbox}
            \begin{bibelbox}{SCHL}{Apg}{17:24-25}
                Der Gott, der die Welt gemacht hat und alles, was darin ist, er, der Herr des Himmels und der Erde, wohnt nicht in Tempeln, die von Menschenhand gemacht sind, und wird auch nicht von Menschenhänden bedient, als ob er etwas nötig hätte, da er selbst allen Leben und Atem und alles gibt.
            \end{bibelbox}
            \item \verbN[b,p]{Er ist gerecht und richtet das Böse.} Ps 7,12; Pred 12,14; Röm 2,6
            \begin{bibelbox}{SCHL}{IMos}{1:1}
                Im Anfang schuf Gott die Himmel und die Erde.
            \end{bibelbox}
            \begin{bibelbox}{SCHL}{Neh}{9:6}
                Du allein bist der HERR. Du hast Himmel und Erde gemacht, das Meer und alles, was darin ist, und du erhältst alles am Leben.
            \end{bibelbox}
            \begin{bibelbox}{SCHL}{Apg}{17:24-25}
                Der Gott, der die Welt gemacht hat und alles, was darin ist, er, der Herr des Himmels und der Erde, wohnt nicht in Tempeln, die von Menschenhand gemacht sind, und wird auch nicht von Menschenhänden bedient, als ob er etwas nötig hätte, da er selbst allen Leben und Atem und alles gibt.
            \end{bibelbox}
            \item \verbN[b,p]{Er ist barmherzig, geduldig und gnädig.} 2Mo 34,6; Ps 103,8; 2Petr 3,9
            \begin{bibelbox}{SCHL}{IMos}{1:1}
                Im Anfang schuf Gott die Himmel und die Erde.
            \end{bibelbox}
            \begin{bibelbox}{SCHL}{Neh}{9:6}
                Du allein bist der HERR. Du hast Himmel und Erde gemacht, das Meer und alles, was darin ist, und du erhältst alles am Leben.
            \end{bibelbox}
            \begin{bibelbox}{SCHL}{Apg}{17:24-25}
                Der Gott, der die Welt gemacht hat und alles, was darin ist, er, der Herr des Himmels und der Erde, wohnt nicht in Tempeln, die von Menschenhand gemacht sind, und wird auch nicht von Menschenhänden bedient, als ob er etwas nötig hätte, da er selbst allen Leben und Atem und alles gibt.
            \end{bibelbox}
            \item \verbN[b,p]{Er ist Liebe und sucht Gemeinschaft mit Menschen.} 1Joh 4,8-10; Jer 31,3; Offb 3,20
            \begin{bibelbox}{SCHL}{IMos}{1:1}
                Im Anfang schuf Gott die Himmel und die Erde.
            \end{bibelbox}
            \begin{bibelbox}{SCHL}{Neh}{9:6}
                Du allein bist der HERR. Du hast Himmel und Erde gemacht, das Meer und alles, was darin ist, und du erhältst alles am Leben.
            \end{bibelbox}
            \begin{bibelbox}{SCHL}{Apg}{17:24-25}
                Der Gott, der die Welt gemacht hat und alles, was darin ist, er, der Herr des Himmels und der Erde, wohnt nicht in Tempeln, die von Menschenhand gemacht sind, und wird auch nicht von Menschenhänden bedient, als ob er etwas nötig hätte, da er selbst allen Leben und Atem und alles gibt.
            \end{bibelbox}
            \item \verbN[b,p]{Er ist treu und hält seine Zusagen.} 4Mo 23,19; 5Mo 7,9; 2Tim 2,13
            \begin{bibelbox}{SCHL}{IMos}{1:1}
                Im Anfang schuf Gott die Himmel und die Erde.
            \end{bibelbox}
            \begin{bibelbox}{SCHL}{Neh}{9:6}
                Du allein bist der HERR. Du hast Himmel und Erde gemacht, das Meer und alles, was darin ist, und du erhältst alles am Leben.
            \end{bibelbox}
            \begin{bibelbox}{SCHL}{Apg}{17:24-25}
                Der Gott, der die Welt gemacht hat und alles, was darin ist, er, der Herr des Himmels und der Erde, wohnt nicht in Tempeln, die von Menschenhand gemacht sind, und wird auch nicht von Menschenhänden bedient, als ob er etwas nötig hätte, da er selbst allen Leben und Atem und alles gibt.
            \end{bibelbox}
            \item \verbN[b,p]{Er ist allwissend und kennt Herz, Gedanken und Zukunft.} Ps 139,1-4; 1Sam 16,7; Jes 46,10
            \begin{bibelbox}{SCHL}{IMos}{1:1}
                Im Anfang schuf Gott die Himmel und die Erde.
            \end{bibelbox}
            \begin{bibelbox}{SCHL}{Neh}{9:6}
                Du allein bist der HERR. Du hast Himmel und Erde gemacht, das Meer und alles, was darin ist, und du erhältst alles am Leben.
            \end{bibelbox}
            \begin{bibelbox}{SCHL}{Apg}{17:24-25}
                Der Gott, der die Welt gemacht hat und alles, was darin ist, er, der Herr des Himmels und der Erde, wohnt nicht in Tempeln, die von Menschenhand gemacht sind, und wird auch nicht von Menschenhänden bedient, als ob er etwas nötig hätte, da er selbst allen Leben und Atem und alles gibt.
            \end{bibelbox}
            \item \verbN[b,p]{Er ist allmächtig und handelt souverän in der Geschichte.} Jer 32,17; Ps 115,3; Dan 4,32
            \begin{bibelbox}{SCHL}{IMos}{1:1}
                Im Anfang schuf Gott die Himmel und die Erde.
            \end{bibelbox}
            \begin{bibelbox}{SCHL}{Neh}{9:6}
                Du allein bist der HERR. Du hast Himmel und Erde gemacht, das Meer und alles, was darin ist, und du erhältst alles am Leben.
            \end{bibelbox}
            \begin{bibelbox}{SCHL}{Apg}{17:24-25}
                Der Gott, der die Welt gemacht hat und alles, was darin ist, er, der Herr des Himmels und der Erde, wohnt nicht in Tempeln, die von Menschenhand gemacht sind, und wird auch nicht von Menschenhänden bedient, als ob er etwas nötig hätte, da er selbst allen Leben und Atem und alles gibt.
            \end{bibelbox}
            \item \verbN[b,p]{Er ist nahe bei den Demütigen, Schwachen und Zerbrochenen.} Ps 34,19; Jes 57,15; Jak 4,6
            \begin{bibelbox}{SCHL}{IMos}{1:1}
                Im Anfang schuf Gott die Himmel und die Erde.
            \end{bibelbox}
            \begin{bibelbox}{SCHL}{Neh}{9:6}
                Du allein bist der HERR. Du hast Himmel und Erde gemacht, das Meer und alles, was darin ist, und du erhältst alles am Leben.
            \end{bibelbox}
            \begin{bibelbox}{SCHL}{Apg}{17:24-25}
                Der Gott, der die Welt gemacht hat und alles, was darin ist, er, der Herr des Himmels und der Erde, wohnt nicht in Tempeln, die von Menschenhand gemacht sind, und wird auch nicht von Menschenhänden bedient, als ob er etwas nötig hätte, da er selbst allen Leben und Atem und alles gibt.
            \end{bibelbox}
            \item \verbN[b,p]{Er ist Retter und Erlöser seines Volkes.} Jes 43,11; Jes 41,14; Tit 2,13-14
            \begin{bibelbox}{SCHL}{IMos}{1:1}
                Im Anfang schuf Gott die Himmel und die Erde.
            \end{bibelbox}
            \begin{bibelbox}{SCHL}{Neh}{9:6}
                Du allein bist der HERR. Du hast Himmel und Erde gemacht, das Meer und alles, was darin ist, und du erhältst alles am Leben.
            \end{bibelbox}
            \begin{bibelbox}{SCHL}{Apg}{17:24-25}
                Der Gott, der die Welt gemacht hat und alles, was darin ist, er, der Herr des Himmels und der Erde, wohnt nicht in Tempeln, die von Menschenhand gemacht sind, und wird auch nicht von Menschenhänden bedient, als ob er etwas nötig hätte, da er selbst allen Leben und Atem und alles gibt.
            \end{bibelbox}
            \item \verbN[b,p]{Er nennt sich Vater für die, die zu ihm gehören.} Mt 6,9; Joh 1,12; 2Kor 6,18
            \begin{bibelbox}{SCHL}{IMos}{1:1}
                Im Anfang schuf Gott die Himmel und die Erde.
            \end{bibelbox}
            \begin{bibelbox}{SCHL}{Neh}{9:6}
                Du allein bist der HERR. Du hast Himmel und Erde gemacht, das Meer und alles, was darin ist, und du erhältst alles am Leben.
            \end{bibelbox}
            \begin{bibelbox}{SCHL}{Apg}{17:24-25}
                Der Gott, der die Welt gemacht hat und alles, was darin ist, er, der Herr des Himmels und der Erde, wohnt nicht in Tempeln, die von Menschenhand gemacht sind, und wird auch nicht von Menschenhänden bedient, als ob er etwas nötig hätte, da er selbst allen Leben und Atem und alles gibt.
            \end{bibelbox}
            \item \verbN[b,p]{Er sagt, dass sein Name geehrt werden soll und dass er keine Konkurrenz durch Götzen duldet.} 2Mo 20,3-5; Jes 42,8; Mt 6,9
            \begin{bibelbox}{SCHL}{IMos}{1:1}
                Im Anfang schuf Gott die Himmel und die Erde.
            \end{bibelbox}
            \begin{bibelbox}{SCHL}{Neh}{9:6}
                Du allein bist der HERR. Du hast Himmel und Erde gemacht, das Meer und alles, was darin ist, und du erhältst alles am Leben.
            \end{bibelbox}
            \begin{bibelbox}{SCHL}{Apg}{17:24-25}
                Der Gott, der die Welt gemacht hat und alles, was darin ist, er, der Herr des Himmels und der Erde, wohnt nicht in Tempeln, die von Menschenhand gemacht sind, und wird auch nicht von Menschenhänden bedient, als ob er etwas nötig hätte, da er selbst allen Leben und Atem und alles gibt.
            \end{bibelbox}
        \end{itemize}
        Die Bibel ist Gottes Wort, welches von Menschen aufgeschrieben wurde. Und das sagt Gott von sich selbst. Nirgends in der Bibel findet ihr etwas, das versucht, Gott zu beweisen. Gott existiert, lebt und handelt - Punkt.
        Aber wie können wir jetzt sicher sein, dass es Gott gibt? Naja, weil wir hier sind. Ohne Gott gibt es auch uns nicht. So einfach ist es. Natürlich ist es nicht so einfach. Wenn es einfach wäre, würden nicht ganze Bibliotheken voll von diesem Thema sein. Ich habe nun vor, mich für die nächste Zeit mit diesem Thema zu beschäftigen. Es geht mir nicht darum, zu beweisen, dass Gott existiert, sondern es geht mir darum, dass unser Vertrauen zu diesem Gott, zu dem wir beten, stärker wird und dass wir ihn anhand der Bibel besser kennenlernen.
        Themen werden sein:
        \begin{enumerate}
            \item Warum der Gott der Bibel und kein anderer?
            \item Gott entdecken
            \item Die Eigenschaften Gottes
            \item Ein einziger Gott, drei göttliche Personen
            \item Gott der Schöpfer
        \end{enumerate}
        Ich habe ein dickes Buch zu Hause in dem diese Themen sehr ausführlich behandelt werden und werde versuchen, euch das in einer Form zu vermitteln, dass auch ich es verstehe.
    \end{block}
    \begin{block}[Ende]
    \end{block}
    \lied[b]{Allein deine Gnade genügt}.
\end{spacing}
\section{Predigt}
\textbf{Nach der Predigt}\newline
Danken für die Predigt.
% \section{Abendmahl}
% \begin{itemize}
%     \item [] Beten für das Brot
%     \item [] Beten für den Wein
% \end{itemize}
\section{Abschluss}
\begin{itemize}
    \item [] \lied[b]{Wenn Friede mit Gott}.
    \item [] \beten[b]{Zum Abschluss}
\end{itemize}
\begin{bibelbox}{SCHL}{IIKor}{13:13}
    Die Gnade des Herrn Jesus Christus und die Liebe Gottes und die Gemeinschaft des Heiligen Geistes sei mit euch allen! Amen
\end{bibelbox}
\end{document}